\documentclass[uplatex, twocolumn,10pt]{jsarticle}

\usepackage[dvipdfmx]{graphicx}
\usepackage{latexsym}
\usepackage{bmpsize}
\usepackage{url}
\usepackage{comment}
\usepackage{layout}
\usepackage[top=30truemm,bottom=30truemm,left=25truemm,right=25truemm]{geometry}

\setlength{\columnsep}{1zw}

\def\Underline{\setbox0\hbox\bgroup\let\\\endUnderline}
\def\endUnderline{\vphantom{y}\egroup\smash{\underline{\box0}}\\}

\newcommand{\ttt}[1]{\texttt{#1}}

\begin{document}

\title{\bf{\LARGE{Sample of \LaTeX  Document} \\ \Large{\LaTeX のサンプルコード}}}
\author{平木場 風太\\宮崎大学大学院}
\date{2019年6月12日(水)}
\maketitle


\section{はじめに}
\begin{figure}[t]
    \begin{center}
        \includegraphics[width=7cm]{image/syokuji_computer.png}
        \caption{パソコンの前でご飯を食べる人のイラスト}
        \label{fig:syokuji_computer}
    \end{center}
\end{figure}

パソコンの前でご飯を食べることはよくある。パソコンの前でご飯を食べる人のイラストを図\ref{fig:syokuji_computer}に示す。
このイラストは、規約の範囲内であれば、個人、法人、商用、非商用問わず無料で利用できることでおなじみの、{\bf かわいいフリー素材 いらすとや}\cite{irasutoya}より引用した。

\section{おわりに}
やっぱり{\bf いらすとや}のイラストはすばらしい。

\begin{thebibliography}{99}
    \bibitem{irasutoya} いらすとや, last access 2019.6.13 \url{https://www.irasutoya.com/}

\section{走れメロス}
メロスは激怒した。必ず、かの邪智暴虐の王を除かなければならぬと決意した。メロスには政治がわからぬ。メロスは、村の牧人である。笛を吹き、羊と遊んで暮して来た。けれども邪悪に対しては、人一倍に敏感であった。きょう未明メロスは村を出発し、野を越え山越え、十里はなれた此のシラクスの市にやって来た。メロスには父も、母も無い。女房も無い。十六の、内気な妹と二人暮しだ。この妹は、村の或る律気な一牧人を、近々、花婿として迎える事になっていた。結婚式も間近かなのである。メロスは、それゆえ、花嫁の衣裳やら祝宴の御馳走やらを買いに、はるばる市にやって来たのだ。先ず、その品々を買い集め、それから都の大路をぶらぶら歩いた。メロスには竹馬の友があった。セリヌンティウスである。今は此のシラクスの市で、石工をしている。その友を、これから訪ねてみるつもりなのだ。久しく逢わなかったのだから、訪ねて行くのが楽しみである。歩いているうちにメロスは、まちの様子を怪しく思った。ひっそりしている。もう既に日も落ちて、まちの暗いのは当りまえだが、けれども、なんだか、夜のせいばかりでは無く、市全体が、やけに寂しい。のんきなメロスも、だんだん不安になって来た。路で逢った若い衆をつかまえて、何かあったのか、二年まえに此の市に来たときは、夜でも皆が歌をうたって、まちは賑やかであった筈だが、と質問した。若い衆は、首を振って答えなかった。しばらく歩いて老爺に逢い、こんどはもっと、語勢を強くして質問した。老爺は答えなかった。メロスは両手で老爺のからだをゆすぶって質問を重ねた。老爺は、あたりをはばかる低声で、わずか答えた。
「王様は、人を殺します。」
「なぜ殺すのだ。」
「悪心を抱いている、というのですが、誰もそんな、悪心を持っては居りませぬ。」
「たくさんの人を殺したのか。」
「はい、はじめは王様の妹婿さまを。それから、御自身のお世嗣を。それから、妹さまを。それから、妹さまの御子さまを。それから、皇后さまを。それから、賢臣のアレキス様を。」
「おどろいた。国王は乱心か。」
「いいえ、乱心ではございませぬ。人を、信ずる事が出来ぬ、というのです。このごろは、臣下の心をも、お疑いになり、少しく派手な暮しをしている者には、人質ひとりずつ差し出すことを命じて居ります。御命令を拒めば十字架にかけられて、殺されます。きょうは、六人殺されました。」
　聞いて、メロスは激怒した。「呆れた王だ。生かして置けぬ。」
　メロスは、単純な男であった。買い物を、背負ったままで、のそのそ王城にはいって行った。たちまち彼は、巡邏の警吏に捕縛された。調べられて、メロスの懐中からは短剣が出て来たので、騒ぎが大きくなってしまった。メロスは、王の前に引き出された。
「この短刀で何をするつもりであったか。言え！」暴君ディオニスは静かに、けれども威厳を以て問いつめた。その王の顔は蒼白で、眉間の皺は、刻み込まれたように深かった。
「市を暴君の手から救うのだ。」とメロスは悪びれずに答えた。
「おまえがか？」王は、憫笑した。「仕方の無いやつじゃ。おまえには、わしの孤独がわからぬ。」
「言うな！」とメロスは、いきり立って反駁した。「人の心を疑うのは、最も恥ずべき悪徳だ。王は、民の忠誠をさえ疑って居られる。」
「疑うのが、正当の心構えなのだと、わしに教えてくれたのは、おまえたちだ。人の心は、あてにならない。人間は、もともと私慾のかたまりさ。信じては、ならぬ。」暴君は落着いて呟き、ほっと溜息をついた。「わしだって、平和を望んでいるのだが。」
「なんの為の平和だ。自分の地位を守る為か。」こんどはメロスが嘲笑した。「罪の無い人を殺して、何が平和だ。」
「だまれ、下賤の者。」王は、さっと顔を挙げて報いた。「口では、どんな清らかな事でも言える。わしには、人の腹綿の奥底が見え透いてならぬ。おまえだって、いまに、磔になってから、泣いて詫びたって聞かぬぞ。」
「ああ、王は悧巧だ。自惚れているがよい。私は、ちゃんと死ぬる覚悟で居るのに。命乞いなど決してしない。ただ、――」と言いかけて、メロスは足もとに視線を落し瞬時ためらい、「ただ、私に情をかけたいつもりなら、処刑までに三日間の日限を与えて下さい。たった一人の妹に、亭主を持たせてやりたいのです。三日のうちに、私は村で結婚式を挙げさせ、必ず、ここへ帰って来ます。」
「ばかな。」と暴君は、嗄れた声で低く笑った。「とんでもない嘘を言うわい。逃がした小鳥が帰って来るというのか。」
「そうです。帰って来るのです。」メロスは必死で言い張った。「私は約束を守ります。私を、三日間だけ許して下さい。妹が、私の帰りを待っているのだ。そんなに私を信じられないならば、よろしい、この市にセリヌンティウスという石工がいます。私の無二の友人だ。あれを、人質としてここに置いて行こう。私が逃げてしまって、三日目の日暮まで、ここに帰って来なかったら、あの友人を絞め殺して下さい。たのむ、そうして下さい。」
　それを聞いて王は、残虐な気持で、そっと北叟笑んだ。生意気なことを言うわい。どうせ帰って来ないにきまっている。この嘘つきに騙された振りして、放してやるのも面白い。そうして身代りの男を、三日目に殺してやるのも気味がいい。人は、これだから信じられぬと、わしは悲しい顔して、その身代りの男を磔刑に処してやるのだ。世の中の、正直者とかいう奴輩にうんと見せつけてやりたいものさ。
「願いを、聞いた。その身代りを呼ぶがよい。三日目には日没までに帰って来い。おくれたら、その身代りを、きっと殺すぞ。ちょっとおくれて来るがいい。おまえの罪は、永遠にゆるしてやろうぞ。」
「なに、何をおっしゃる。」
\end{thebibliography}

\layout

\end{document}
